\section{빛의 섬광으로 주는 도파민}
\label{sec:dishdopamine}

배양에 도파민을 교사로 준 직접 실험으로 우리가 아는 것은 하나뿐이다. 연구자들이 원격으로 접속하는 오가노이드 플랫폼에서 수행되었다~\cite{Jordan2024}. 오가노이드를 적시는 액체에 빛에 반응하는 “케이지”에 갇힌 도파민을 넣었다. 묶여 있는 분자는 수용체에 작용하지 않는다. 케이지 도파민의 농도는 $1$~mM이다. 네트워크에 보상을 주어야 할 때 자외선을 비추면 케이지가 분해되고 도파민이 부피 속으로 풀려난다.

루프는 이렇게 짜여 있다. 같은 자극을 되풀이해서 넣고, 그 자극이 이전보다 많은 스파이크를 일으키면 섬광을 켜서 도파민을 방출한다. $240$번 반복하는 동안 유발 스파이크 수는 전체적으로 늘었다. 저자들 스스로 이런 실험 하나만으로는 증가가 도파민 때문이라고 결론 내릴 수 없다고 쓴다~\cite{Jordan2024}.

엔지니어에게 이것은 3요소 규칙~\eqref{eq:threefactor}을 접시 안에 조립하려는 첫 시도이다. 송신자와 수신자의 활동이 흔적을 남기고, 공통 신호가 변화를 굳힐지 결정한다. 그러나 여기서 신호는 양수뿐이다. 보상이 있거나 없거나이며, 장치는 음의 오차 $\delta<0$를 내지 못한다. 게다가 도파민은 \eqref{eq:lowpass}의 농도처럼 퍼지고 천천히 제거되므로, 섬광이 잦으면 그저 전체 수준만 올라갈 수 있다. 학습을 이것과 구별하려면 대조 실험 두 가지가 필요하다. 네트워크의 응답과 무관한 무작위 시점에 같은 수의 섬광을 주는 것, 그리고 케이지 도파민 없이 같은 섬광을 주는 것이다. 또 휴식 뒤의 확인도 필요하다. 도파민이 사라진 뒤에도 변화가 남았는가?

\section{도파민이 실리콘을 가르친다}
\label{sec:biohybrid}

반대 방향의 실험에서는 살아 있는 세포가 전자 소자를 가르친다~\cite{Keene2020}. 도파민을 분비하는 세포를 유기 뉴로모픽 소자 위의 미세 채널에서 키운다. 이 소자는 트랜지스터이며, 채널의 전도도가 시냅스 가중치 역할을 한다. 세포에서 나온 도파민이 소자의 전극에서 산화되면서 전도도를 바꾼다. 이 장치는 틈새에서 신경전달물질을 치우는 메커니즘도 재현하므로, 가중치를 오래 바꿀 수 있을 뿐 아니라 원래대로 되돌릴 수도 있다.

회로로 보면 이것은 화학 입력을 받는 누설 적분기이다. 가중치는 도파민 흐름을 누적하고, “제거”가 얼마나 빨리 잊는지를 정한다. \eqref{eq:lowpass}와 같은 RC 회로이다. 중요한 것은 연결 방향이다. \ref{sec:debugging}장의 프로브는 브로드캐스트 레지스터를 보지 못한다. 그것이 화학적이기 때문이다. 이 소자는 바로 화학 레지스터를 읽어서 전기적 가중치로 바꾼다. 뇌-컴퓨터 인터페이스에게 이것은 데이터 버스뿐 아니라 제어 레지스터까지 듣는 센서로 가는 길이다.

\section{자기 도파민 뉴런을 가진 오가노이드}
\label{sec:midbrainorg}

인간 줄기세포로 \nt{DOPA} 뉴런이 있는 뇌 영역을 닮은 오가노이드를 키울 수 있다. 그 안에는 도파민을 분비하는, 작동하는 도파민 뉴런이 있다~\cite{Jo2016}. 이런 오가노이드는 주로 질병 연구를 위해 키운다. 그 도파민이 다른 네트워크의 교사 역할을 한 실험은 찾지 못했다.

살아 있는 뉴런으로 만든 기계에게 이것은 빠진 부품이다. \ref{sec:livingmachine}장의 테스트 벤치는 제어 레지스터 없는 데이터 버스와 흐름 제어이다. 여기에는 브로드캐스트 신호의 원천이 이미 있다. 모자란 것은 비평가, 즉 예측 오차~\eqref{eq:ctd}를 계산하고 이 뉴런들을 조종하는 네트워크이다. 피질 오가노이드를 이런 오가노이드와 연결해서 오차에 따라 도파민이 분비되게 하는 것은 열린 문제이며, 아직은 실험이 아니라 가설이다.

\section{도파민 없이 막대의 균형을 잡는 오가노이드}
\label{sec:cartpole}

최근 실험 가운데 가장 완전한 것은 도파민을 전혀 쓰지 않는다~\cite{Robbins2026}. 생쥐 피질 오가노이드를 “카트 위의 막대” 과제에 닫힌 루프로 연결했다. 오가노이드의 활동이 카트를 움직이고, 막대의 위치는 자극으로 되돌아 들어간다. 시도와 시도 사이에 오가노이드는 짧은 고주파 학습 자극을 받는다. 어떤 자극을 줄지는 강화 학습 프로그램이 고른다.

측정된 결과는 다음과 같다.
\begin{itemize}
\item 대부분의 오가노이드에서 프로그램이 고른 학습 신호가 무작위로 고른 신호나 신호가 없을 때보다 좋은 결과를 냈다.
\item $45$분 휴식 뒤에는 향상이 남지 않았다.
\item 두 종류의 글루탐산 수용체 AMPA와 NMDA를 모두 차단하면 향상이 사라졌다.
\end{itemize}

마지막 항목은 학습된 것이 어디에 사는지 말해 준다. \nt{GLUT} 시냅스 안이며, \ref{sec:weightstore}절에서 입력과 출력의 동시성 검출기로 작동하는 수용체를 통해서이다. 두 번째 항목은 무엇이 모자란지 보여 준다. 빠른 변화는 있지만 굳히기는 없다. \ref{sec:motorlearning}장에서 느린 학습자 없이 빠른 학습자만 있을 때와 같다. 교사의 역할도 바뀌었다. 전류 패턴을 고르는 엔지니어 대신 프로그램이 들어섰지만, 교사는 여전히 접시 바깥에 있다.

배양을 전극 배열에서 학습시킨 더 이른 실험들 중에서도 도파민을 쓴 것은 하나도 찾지 못했다. 학습 신호는 모두 전기 자극이었다.

\section{누가 배양을 가르치는가}

\begin{table}[t]
\centering
\footnotesize
\setlength{\tabcolsep}{4pt}
\renewcommand{\arraystretch}{1.15}
\begin{tabular}{>{\raggedright}p{2.7cm}>{\raggedright}p{2.5cm}>{\raggedright}p{3.2cm}>{\raggedright\arraybackslash}p{4.4cm}}
\toprule
실험 & 가르치는 주체 & 학습 신호 & 알려진 것 \\
\midrule
원하는 응답에서 자극 멈추기 & 엔지니어 & 자극의 종료 & 응답이 굳어진다 --- 측정됨 \\
DishBrain (\ref{sec:dishbrain}장) & 엔지니어 & 예측 가능한 전류 또는 무작위 전류 & 세션 안의 유의미한 랠리 증가는 완전한 되먹임일 때만, 침묵일 때는 효과가 더 작다 --- 측정됨. 세션 간에는 불안정. 원인은 논쟁 중 \\
카트 위의 막대 & 강화 학습 프로그램 & 프로그램이 고른 고주파 자극 & 무작위보다 낫지만 휴식 뒤에는 남지 않는다 --- 측정됨 \\
섬광으로 주는 도파민 & “스파이크가 늘면 보상” 규칙 & 빛으로 풀려난 도파민 & 스파이크 증가 --- 관찰. 도파민의 역할은 증명되지 않음 \\
바이오하이브리드 시냅스 & 살아 있는 세포 & 세포에서 나온 도파민 & 소자 가중치의 장기 변화와 복귀 --- 측정됨 \\
\nt{DOPA} 뉴런을 가진 오가노이드 & --- & --- & 도파민 원천은 있다. 교사로 쓰인 적은 없다 \\
\bottomrule
\end{tabular}
\caption{누가 무엇으로 칩 위의 살아 있는 뉴런을 가르치는가.}
\label{tab:dishteachers}
\end{table}

배양 자체가 학습하는 모든 실험에서 교사는 아직 바깥에 있다(표~\ref{tab:dishteachers}). 엔지니어, 프로그램, 또는 엔지니어가 정한 규칙이다. 도파민이 접시 안에 들어온 것은 단 한 번뿐이고, 그 역할은 아직 증명해야 한다. \ref{sec:dofa}장에서처럼 비평가, 오차, 흔적을 갖춘 진짜 브로드캐스트 채널을 접시 안에 조립하는 것은 아직 아무도 내딛지 않은 다음 걸음이다.
